\section{Research programme through alpha: exact pairing and analytic obligations}
\label{sec:alpha-programme}

This independently readable chapter states the research programme, proves
its intermediate identities, and compares its remaining obligations with
established approaches. The comparison concerns the cited works, not every
proposed RH proof. An equivalence theorem, a conditional theorem, and a
speculative model are not all failed mathematics.

The active route is exact cancellation between sign-specific ordinal pairs,
retaining unfinished pairs at every integer cutoff. No estimate of RH
strength is assumed. The prime number theorem (PNT) and elementary squarefree
counting are used explicitly as unconditional background. Statements are
separated into proved results, classical transfers, unresolved targets,
and excluded universal mechanisms.

\subsection{Step 0: independent notation and target}
\label{subsec:a12-definitions}
Let $\mu$ be the usual M\"obius function, $p_j$ the $j$th prime, and $p_0=0$.
For real $x\ge1$ define the step functions
\begin{equation}
\label{eq:a12-counts}
N_\pm(x)=\#\{n\le x:\mu(n)=\pm1\},\quad
Q(x)=N_+(x)+N_-(x),\quad M(x)=N_+(x)-N_-(x).
\end{equation}
The transform and its cumulative sum are
\begin{equation}
\label{eq:a12-S}
S(n)=
\begin{cases}
p_{N_+(n)},&\mu(n)=1,\\
-p_{N_-(n)},&\mu(n)=-1,\\
1,&\mu(n)=0,
\end{cases}
\qquad SS(x)=\sum_{n\le x}S(n).
\end{equation}
Set $w_j=p_j-j$ and $D_w(m)=\sum_{j=1}^m w_j$, with $D_w(0)=0$.
For $\mu(n)\ne0$, let $j(n)$ be the occurrence number within its own sign.
The index $j(n)$ is not used at nonsquarefree integers. Define
\begin{equation}
\label{eq:a12-series}
\alpha(s)=\sum_{n\ge1}S(n)n^{-s},\qquad
\delta(s)=\sum_{\substack{n\ge1\\\mu(n)\ne0}}\mu(n)w_{j(n)}n^{-s}.
\end{equation}
Let $s_1<s_2<\cdots$ enumerate the squarefree integers. Write
\begin{equation}
\label{eq:a12-pairs}
\begin{aligned}
\varepsilon_q&=\mu(s_q),&
R(q)&=\sum_{\ell\le q}\varepsilon_\ell=M(s_q),\qquad R(0)=0,\\
\alpha_j&=\text{ordinal position of the $j$th $+1$},&
\beta_j&=\text{ordinal position of the $j$th $-1$},\\
a_j&=s_{\alpha_j},&b_j&=s_{\beta_j},\\
d_j&=\beta_j-\alpha_j,&\Delta_j&=b_j-a_j.
\end{aligned}
\end{equation}
The function $\alpha(s)$ differs from the ordinal $\alpha_j$.
The physical positions $a_j,b_j$ avoid a clash with the analytic blocks
$A(s),B(s)$. Never identify $d_j$, $\Delta_j$, and $w_j$.
The notation $D_w$ also differs from a lifetime generating function $D(z)$.
Finally set
\begin{equation}
\label{eq:a12-H}
c_0=\frac3{\pi^2},\qquad
\Theta=\sup_{\rho\ {\rm nontrivial}}\Re\rho,\qquad
\mathcal H=\{s:\Re s>3/2\}.
\end{equation}
The target is exclusion of nontrivial zeros with $\Re\rho>1/2$.
The functional equation supplies the other half of the strip.

\subsection{Step 1: decomposition and genuine pole statements}
\label{subsec:a12-decomposition}

\begin{proposition}[Exact decomposition and domains]
\label{prop:a12-decomposition}
The series in \eqref{eq:a12-series} converge absolutely for $\Re s>2$,
and in that domain
\begin{equation}
\label{eq:a12-ABCD}
\boxed{\alpha(s)=A(s)+B(s)+C(s)+\delta(s),}
\end{equation}
where
\begin{align}
A(s)&=\zeta(s)-\frac{\zeta(s)}{\zeta(2s)},\label{eq:a12-A}\\
B(s)&=c_0\frac{s}{(s-1)\zeta(s-1)},\label{eq:a12-B}\\
C(s)&=\frac{s}{2}\int_1^\infty
M(x)(Q(x)+1-2c_0x)x^{-s-1}\,dx.\label{eq:a12-C}
\end{align}
The blocks $A,B$ are meromorphic on $\mathbb C$.
The integral for $C$ is holomorphic on $\mathcal H$ unconditionally;
more precisely, its domain includes $\Re s>\Theta+1/2$.
\end{proposition}
\begin{proof}
The bound $p_j\ll j\log(2j)$ majorises both coefficient sequences by
$O(n\log(2n))$. The signed ordinal coefficients have the partial sum
\begin{equation}
\label{eq:a12-ordinal}
\sum_{\substack{n\le x\\\mu(n)\ne0}}\mu(n)j(n)
=\frac{N_+(N_++1)-N_-(N_-+1)}2=\frac{M(x)(Q(x)+1)}2.
\end{equation}
Partial summation gives the integral of this expression against
$sx^{-s-1}\,dx$. Splitting $Q+1=2c_0x+(Q+1-2c_0x)$ and using
\begin{equation}
\label{eq:a12-Mellin}
\int_1^\infty M(x)x^{-z-1}\,dx=\frac1{z\zeta(z)}
\qquad(\Re z>1)
\end{equation}
gives $B+C$. The zero row contributes $A$.
The unconditional estimates $Q(x)=2c_0x+O(\sqrt x)$ and $|M(x)|\le x$
give local uniform absolute convergence of the $C$ integral for $\Re s>3/2$.
The classical bound $M(x)=O_\epsilon(x^{\Theta+\epsilon})$ gives the
refined domain \cite{Titchmarsh1986}.
\end{proof}

\begin{proposition}[Poles of the ordinal main block]
\label{prop:a12-poles}
A nontrivial zero $\rho$ of $\zeta$ gives a pole of $B$ at $1+\rho$ with
the same order. For a simple zero,
\begin{equation}
\label{eq:a12-residue}
\operatorname*{Res}_{s=1+\rho}B(s)
=c_0\frac{1+\rho}{\rho\zeta'(\rho)}.
\end{equation}
There is also a simple pole at $s=1$ of residue $-2c_0$, and the shifted
trivial-zero poles $s=1-2k$, $k\ge1$. The point $s=2$ is a zero, not a pole.
\end{proposition}
\begin{proof}
Use the local factorisation of $\zeta(s-1)$ at its zeros.
The prefactor $c_0s/(s-1)$ is finite and nonzero at every shifted zero.
At $s=1$ use $\zeta(0)=-1/2$; at $s=2$ use the pole of $\zeta$ at $1$.
\end{proof}

No corresponding pole count for the full $\alpha$ has been proved.
Critical-line zero theorems transfer to $B$, not automatically to $\alpha$:
continuation and noncancellation of the other blocks must first be checked.
Holomorphy of $\delta$ alone does not remove a pole of $B$ from $\alpha$.
Nor does extending $\delta$ across $3/2$ extend $C$ across that line.

\subsection{Step 2: exact pairing without linearisation}
\label{subsec:a12-exact-pairing}
\begin{proposition}[Pairs, oriented integrals, and unit cuts]
\label{prop:a12-exact-pairing}
For $\Re s>2$,
\begin{equation}
\label{eq:a12-delta-pairs}
\delta(s)=\sum_{j\ge1}w_j(a_j^{-s}-b_j^{-s}).
\end{equation}
For each pair and every complex $s$,
\begin{equation}
\label{eq:a12-exact-integral}
a_j^{-s}-b_j^{-s}=s\int_{a_j}^{b_j}x^{-s-1}\,dx.
\end{equation}
With
\begin{equation}
\label{eq:a12-U}
u_q=\frac{q+R(q)}2,\quad v_q=\frac{q-R(q)}2,\quad
U(q)=D_w(u_q)-D_w(v_q),
\end{equation}
the initial half-plane has the representation
\begin{equation}
\label{eq:a12-unit-cut}
\boxed{\delta(s)=\sum_{q\ge1}U(q)(s_q^{-s}-s_{q+1}^{-s}).}
\end{equation}
\end{proposition}
\begin{proof}
Absolute convergence allows occurrence matching.
Differentiating $x^{-s}$ proves the oriented integral, including $s=0$.
Telescoping gives
\begin{equation}
\label{eq:a12-telescope}
a_j^{-s}-b_j^{-s}=\operatorname{sgn}(d_j)
\sum_{q=\min(\alpha_j,\beta_j)}^{\max(\alpha_j,\beta_j)-1}
(s_q^{-s}-s_{q+1}^{-s}).
\end{equation}
At cut $q$ the open labels are exactly
$\min(u_q,v_q)<j\le\max(u_q,v_q)$, all with orientation
$\operatorname{sgn}R(q)$. Their signed weight is $U(q)$.
The bound $|U(Q(x))|\ll x^2\log(2x)$ gives an integrable absolute majorant
$|s|x^{1-\Re s}\log(2x)$ when $\Re s>2$, justifying the rearrangement.
\end{proof}
In particular,
\begin{equation}
\label{eq:a12-U-block}
U(q)=\operatorname{sgn}R(q)
\sum_{j=\min(u_q,v_q)+1}^{\max(u_q,v_q)}w_j.
\end{equation}
The kernel in \eqref{eq:a12-unit-cut} is local, but its coefficient still
contains every long open pair. At one cut all these contributions have the
same sign. Useful cancellation must occur across different cuts.

PNT already gives $\alpha_j,\beta_j\sim2j$ and $a_j,b_j\sim j/c_0$.
Thus $d_j=o(j)$ and $\Delta_j=o(j)$ unconditionally: eventually excluding
a macroscopic relative separation does not require RH. This relative
smallness supplies no square-root cancellation, and a Taylor argument
would still need uniform summable errors. The exact formulas avoid that
requirement entirely.

\subsection{Step 3: the nonpolynomial duration observable}
\label{subsec:a12-duration}
For this step take a balanced finite word of length $Q=2J$, so $R(Q)=0$.
All its $J$ pairs are complete. Define half-open cut intervals and duration tails:
\begin{equation}
\label{eq:a12-halfopen}
I_j=[\min(\alpha_j,\beta_j),\max(\alpha_j,\beta_j)),\qquad
V_J(r)=\sum_{j=1}^J(|d_j|-r)_+.
\end{equation}
\begin{proposition}[Two-cut identity]
\label{prop:a12-two-cut}
For $0\le r<Q$,
\begin{equation}
\label{eq:a12-V}
\boxed{V_J(r)=\frac12\sum_{q=0}^{Q-r-1}
(|R(q)+R(q+r)|-r)_+.}
\end{equation}
For $r\ge Q$ use an empty sum. In particular,
\begin{equation}
\label{eq:a12-V-zero}
V_J(0)=\sum_{q=0}^{Q-1}|R(q)|\ge J.
\end{equation}
For $0<z<1$,
\begin{equation}
\label{eq:a12-Phi}
\Phi_J(z):=\sum_{j=1}^Jz^{|d_j|}
=J-(1-z)V_J(0)+(1-z)^2\sum_{r=1}^{Q-1}z^{r-1}V_J(r).
\end{equation}
\end{proposition}
\begin{proof}
Using $N_\pm(q)=(q\pm R(q))/2$ in the ordinal clock, pairs opening with
plus and crossing both cuts have labels $N_-(q+r)<j\le N_+(q)$.
Those opening with minus have labels $N_+(q+r)<j\le N_-(q)$.
Their disjoint counts give
\begin{equation}
\label{eq:a12-Cr}
C_r(q)=
(N_+(q)-N_-(q+r))_++(N_-(q)-N_+(q+r))_+
=\tfrac12(|R(q)+R(q+r)|-r)_+.
\end{equation}
An interval of length $\ell$ contains both cuts for $(\ell-r)_+$ choices
of $q$. Double counting proves \eqref{eq:a12-V}. The case $r=0$ and
$|d_j|\ge1$ give \eqref{eq:a12-V-zero}.
Finally, sum the finite identity
\begin{equation}
\label{eq:a12-single-duration}
z^\ell=1-(1-z)\ell+
(1-z)^2\sum_{r=1}^{\ell-1}z^{r-1}(\ell-r)
\end{equation}
over the pairs.
\end{proof}

\begin{proposition}[Order survives a one-time height histogram]
\label{prop:a12-order-example}
The words $\mathrm{I}=++-+--+-$ and $\mathrm{II}=++--++--$ have the same
height histogram but different duration tails and transforms.
\end{proposition}
\begin{proof}
Their height lists are $(1,2,1,2,1,0,1,0)$ and $(1,2,1,0,1,2,1,0)$:
both contain two zeros, four ones, and two twos.
Their durations are $(2,3,2,1)$ and $(2,2,2,2)$, respectively. Therefore
\begin{equation}
\label{eq:a12-order-separation}
V_{\rm I}(2)=1,\quad V_{\rm II}(2)=0,\qquad
\Phi_{\rm I}(z)-\Phi_{\rm II}(z)=z(1-z)^2>0.
\end{equation}
Their quadratic duration sums also differ, $18$ versus $16$.
Not every quadratic functional is insensitive to order.
\end{proof}

Order sensitivity is not yet arithmetic specificity. If $F_+,F_-$ are the
distribution functions of the empirical probability measures
$J^{-1}\sum_j\delta_{\alpha_j}$ and $J^{-1}\sum_j\delta_{\beta_j}$, then
\begin{equation}
\label{eq:a12-Wasserstein}
\frac{V_J(0)}J=\frac1J\sum_j|\alpha_j-\beta_j|
=\int_{\mathbb R}|F_+(t)-F_-(t)|\,dt.
\end{equation}
This is the classical one-dimensional monotone Wasserstein transport
identity \cite{Vallender1974}. Neither monotone matching nor its first
absolute cost is claimed as new.
\subsubsection{What fixed-lag arithmetic input actually controls}
\label{subsubsec:a12-local-input}
The squarefree-clock two-point theorem obtained from Tao's logarithmic
two-point theorem \cite{Tao2016} gives, for every fixed $r\ge1$,
\begin{equation}
\label{eq:a12-window-input}
\frac1{\log X}\sum_{q\le X}\frac{(R(q+r)-R(q))^2}{q}\longrightarrow r.
\end{equation}
The clock transfer keeps a physical separation and a squarefree-sieve
modulus fixed when applying the two-affine-form theorem.
The error from omitted prime squares tends to zero as the sieve cutoff
subsequently tends to infinity. Large physical separations have vanishing
logarithmic tail because
$\sum_{q\le X}(s_{q+r}-s_q)=O_r(X)$ by telescoping.
Expanding the increment square then gives $r$ diagonal terms and finitely
many correlations. These successive limits do not give uniformity in
growing $r$.

The precise relation to the duration observable is
\begin{equation}
\label{eq:a12-local-error}
\left|C_r(q)-\left(|R(q)|-\frac r2\right)_+\right|
\le\frac12|R(q+r)-R(q)|.
\end{equation}
It follows from the $1$-Lipschitz property of absolute value and positive part.
Cauchy--Schwarz gives
\begin{equation}
\label{eq:a12-local-average}
\limsup_{X\to\infty}\frac1{\log X}
\sum_{q\le X}\frac{|C_r(q)-(|R(q)|-r/2)_+|}{q}
\le\frac{\sqrt r}{2}.
\end{equation}
The unestimated main term depends on the absolute height. A theorem on
fixed increments is therefore not automatically a bound for $V_J(r)$,
even for fixed $r$. Moreover $V_J$ is unsigned and lacks the weights $w_j$;
both transfers to $\delta$ require additional argument.

\begin{remark}[Height, area, and prime energy are different]
\label{rem:a12-scales}
Let $H=\max_{q\le Q}|R(q)|$. The $H$ pairs active at a cut of height
$\pm H$ each have duration at least $H$: for the $t$th active label,
at least $H-t$ steps separate its opening from the cut and at least $t$
steps separate the cut from its closing. Hence
\begin{equation}
\label{eq:a12-V-height}
V_J(r)\ge H(H-r)_+.
\end{equation}
In particular $V_J(0)$ is an area, not a height.
The bound $V_J(0)\ll_\epsilon J^{1/2+\epsilon}$ for every $\epsilon>0$
is impossible because $V_J(0)\ge J$.
A square-root-plus-$\epsilon$ height bound instead gives an area bound
$O_\epsilon(J^{3/2+\epsilon})$.
The cubic negative-excursion prime-separation energy is a different functional.
\end{remark}

\subsection{Step 4: holomorphy for a fixed duration cutoff}
\label{subsec:a12-short-pairs}

\begin{theorem}[Unconditional bounded-duration holomorphy]
\label{thm:a12-short}
For every fixed $L\ge1$ the series
\begin{equation}
\label{eq:a12-short}
\delta_{\le L}(s)=\sum_{|d_j|\le L}w_j(a_j^{-s}-b_j^{-s})
\end{equation}
converges normally by pairs in $\Re s>1$ and is holomorphic there.
For $\sigma=\Re s>1$,
\begin{equation}
\label{eq:a12-short-bound}
\sum_{|d_j|\le L}w_j|a_j^{-s}-b_j^{-s}|
\ll L|s|\left(\frac{\log2}{\sigma-1}+\frac1{(\sigma-1)^2}\right).
\end{equation}
The proof uses no Mertens estimate, density theorem, or linearisation.
\end{theorem}
\begin{proof}
At physical cut $x$, put $q=Q(x)$. An active ordinal interval of length
at most $L$ opened within the last $L$ positions. There are at most $L$
such intervals, and each has label $j\le q\le x$, so
$w_j\ll x\log(2x)$. The exact integral and Tonelli give
\begin{align}
\sum_{|d_j|\le L}w_j|a_j^{-s}-b_j^{-s}|
&\le |s|\int_1^\infty x^{-\sigma-1}
\sum_{\substack{|d_j|\le L\\\min(a_j,b_j)\le x<\max(a_j,b_j)}}w_j\,dx
\label{eq:a12-short-majorant}\\
&\ll L|s|\int_1^\infty x^{-\sigma}\log(2x)\,dx.
\nonumber
\end{align}
The last integral equals the parenthesis in \eqref{eq:a12-short-bound}.
On any compact set use its minimum real part and maximum $|s|$ to obtain
a summable majorant for the suprema of the paired terms.
\end{proof}

This theorem also holds for arbitrary sign-occurrence lists in an
increasing integer clock and weights $|w_j|\ll j\log(2j)$.
Its generality is useful, but it is not an arithmetic cancellation theorem.

For $\Re s>2$, dominated convergence gives
$\delta_{\le L}\to\delta$ locally uniformly as $L\to\infty$.
The displayed bound grows with $L$; it does not justify this limit on
$\mathcal H$.

\subsection{Step 5: long pairs and the sign involution}
\label{subsec:a12-long-pairs}
Large rank and large duration are different: an alternating word has
$|d_j|=1$ at every rank, however large $w_j$ becomes.

\begin{proposition}[Universal antisymmetry does not delete the long tail]
\label{prop:a12-no-universal-tail}
Sign reversal and ordinal matching alone do not justify discarding all
pairs of duration greater than a fixed $L$.
\end{proposition}
\begin{proof}
For $H>L$, take the auxiliary balanced word $(+1)^H(-1)^H$ on any
increasing positive clock. Every duration equals $H$, but for real
$\sigma>0$,
\begin{equation}
\label{eq:a12-tail-counterexample}
\delta_{\le L}(\sigma)=0,\qquad
\sum_{j=1}^H w_j(s_j^{-\sigma}-s_{H+j}^{-\sigma})>0.
\end{equation}
There are no short pairs. This disproves the claimed universal deletion,
not a possible cancellation specific to the genuine M\"obius word.
\end{proof}

The global sign flip gives
\begin{equation}
\label{eq:a12-signflip}
R\mapsto-R,\qquad d_j\mapsto-d_j,\qquad \delta(s)\mapsto-\delta(s).
\end{equation}
It has no action on $s$ and gives no equation under $s\mapsto3-s$.
The zero row of $\alpha$ is unchanged, so the whole transform does not
simply change sign either.

If a universal functional satisfies
$F(\varepsilon)\ge c\,\Psi(\|d(\varepsilon)\|_\infty)$ for both signs,
then its evenisation
$[F(\varepsilon)+F(-\varepsilon)]/2$ has the same lower bound.
Thus odd cancellation alone cannot explain such universal coercivity.
This does not exclude all even, nonpolynomial, or arithmetic-specific
functionals.

\subsection{Step 6: the precise original-order target}
\label{subsec:a12-cauchy}
For integer $X$ let
\begin{equation}
\label{eq:a12-delta-X}
\delta_X(s)=\sum_{j\ge1}w_j
\bigl(\mathbf1_{a_j\le X}a_j^{-s}-\mathbf1_{b_j\le X}b_j^{-s}\bigr).
\end{equation}
This is a finite sum in the original integer ordering.

\subsubsection{The unfinished-pair term}
\label{subsubsec:a12-boundary}
Put $m_X=\min(N_+(X),N_-(X))$ and $h_X=|N_+(X)-N_-(X)|$. Exactly,
\begin{equation}
\label{eq:a12-boundary-split}
\delta_X(s)=\sum_{j\le m_X}w_j(a_j^{-s}-b_j^{-s})+E_X(s),
\end{equation}
where
\begin{equation}
\label{eq:a12-boundary-term}
E_X(s)=
\begin{cases}
\displaystyle\sum_{j=m_X+1}^{m_X+h_X}w_j a_j^{-s},
 &N_+(X)>N_-(X),\\[2mm]
\displaystyle-\sum_{j=m_X+1}^{m_X+h_X}w_j b_j^{-s},
 &N_-(X)>N_+(X),\\
0,&N_+(X)=N_-(X).
\end{cases}
\end{equation}
At a real argument all terms of $E_X$ have the same sign.
Convergence along balanced prefixes or after completing every pair does
not justify convergence at arbitrary integer cutoffs.

\subsubsection{Unresolved Cauchy estimate}
\label{subsubsec:a12-target}
The central target is, for every compact $K\subset\mathcal H$,
\begin{equation}
\label{eq:a12-target}
\boxed{\lim_{X\to\infty}
\sup_{\substack{Y\in\mathbb N\\Y\ge X}}
\sup_{s\in K}|\delta_Y(s)-\delta_X(s)|=0.}
\end{equation}
A concrete sufficient estimate is
$|\delta_Y(s)-\delta_X(s)|\le\eta_K(X)$ for all such $Y,s$, where an
independently controlled $\eta_K(X)$ tends to zero.
Complete short pairs, complete long pairs, and pairs cut by either endpoint
must be handled with compatible uniformity.

This is locally uniform convergence of original partial sums, not absolute
convergence of individual integer terms, whose threshold remains $2$.
A different grouping that gives analytic continuation is not automatically
a proof of \eqref{eq:a12-target}.

\subsection{Step 7: the rigorous transfer to RH}
\label{subsec:a12-transfer}
\begin{proposition}[The exact coefficient sum and its scale]
\label{prop:a12-W-scale}
Let $W(x)=\sum_{n\le x,\ \mu(n)\ne0}\mu(n)w_{j(n)}$. Then
\begin{equation}
\label{eq:a12-W}
W(x)=D_w(N_+(x))-D_w(N_-(x))
=SS(x)-\lfloor x\rfloor+Q(x)-\frac{M(x)(Q(x)+1)}2.
\end{equation}
There is $e(x)\to0$ such that, for sufficiently large $x$,
\begin{equation}
\label{eq:a12-W-scale}
W(x)=c_0x\log x\,M(x)(1+e(x)),
\end{equation}
with both sides zero when $M(x)=0$.
\end{proposition}
\begin{proof}
The difference of cumulative weights gives the first identity, and
\eqref{eq:a12-ordinal} gives the second. PNT implies $Q(x)\sim2c_0x$ and
$M(x)=o(x)$, hence both counters are $c_0x(1+o(1))$.
Uniformly between them, $w_j=p_j-j\sim j\log j\sim c_0x\log x$.
The block has signed length $M(x)$. This uses PNT, not an RH-sized bound.
\end{proof}

\begin{theorem}[Original-order criterion and RH]
\label{thm:a12-transfer}
The Cauchy target \eqref{eq:a12-target} is equivalent to RH.
It is also equivalent to convergence of the ordinary series $\delta(\sigma)$
for every real $\sigma>3/2$.
\end{theorem}
\begin{proof}
Real convergence and Kronecker's lemma give $W(x)=o(x^\sigma)$.
Consequently \eqref{eq:a12-W-scale} gives
\begin{equation}
\label{eq:a12-Kronecker}
M(x)=o\left(\frac{x^{\sigma-1}}{\log x}\right).
\end{equation}
Taking $\sigma=3/2+\epsilon/2$ for each $\epsilon>0$ yields the
Littlewood criterion for RH \cite{Littlewood1912,Titchmarsh1986}.
Its needed direction follows because these estimates make
$z\int_1^\infty M(x)x^{-z-1}\,dx$ holomorphic for $\Re z>1/2$.
It agrees with $1/\zeta(z)$ for $\Re z>1$, so no zero is possible to
the right of $1/2$; symmetry completes RH.

Conversely, under RH, $M(x)=O_\eta(x^{1/2+\eta})$ gives
$W(x)=O_\eta(x^{3/2+\eta}\log x)$.
For compact $K\subset\mathcal H$ choose
$0<\eta<\min_{s\in K}(\Re s-3/2)$. Partial summation gives exactly
\begin{equation}
\label{eq:a12-Abel-tail}
\delta_Y(s)-\delta_X(s)
=W(Y)Y^{-s}-W(X)X^{-s}
+s\int_X^Y W(x)x^{-s-1}\,dx.
\end{equation}
All terms tend to zero uniformly in $Y\ge X$ and $s\in K$.
\end{proof}

This is a classical transfer in new coefficients, not a proof of its
unresolved premise. In the forward direction the Mertens estimate is
a conclusion. Assuming it while proving the premise would be circular.

The standard abscissa formula and the classical identification of the
Mertens growth exponent with $\Theta$ also yield
\begin{equation}
\label{eq:a12-abscissae}
\sigma_c(\alpha)=\sigma_c(\delta)=1+\Theta.
\end{equation}
Indeed $W$ has growth exponent $1+\Theta$; the signed prime part of $SS$
has the same exponent, whereas $\lfloor x\rfloor-Q(x)=O(x)$.
The abscissae refer to the original ordering, not to arbitrary paired
rearrangements. An abscissa is not automatically a line of poles.

The ordinary series $\alpha(3/2)$ diverges. Otherwise Kronecker's lemma
would give $SS(x)=o(x^{3/2})$ and the same signed prime-block asymptotic
would force $M(x)=o(\sqrt x/\log x)$, contrary to
$M(x)=\Omega_\pm(\sqrt x)$ \cite{Ng2004,OdlyzkoTeRiele1985}.
The stronger constants of the Mertens-conjecture disproof are unnecessary.
Divergence of this series does not decide whether a continuation could
have a finite value at the real point $3/2$.
\subsection{Step 8: the separate continuation-and-poles route}
\label{subsec:a12-continuation}
The meromorphic route has a different sufficient objective:
\begin{equation}
\label{eq:a12-analytic-criterion}
\boxed{
\begin{aligned}
\mathrm{RH}&\Longleftrightarrow B\in\mathcal O(\mathcal H)\\
&\Longleftrightarrow
\alpha-\delta\text{ admits holomorphic continuation to }\mathcal H.
\end{aligned}}
\end{equation}
The first equivalence follows from Proposition~\ref{prop:a12-poles} and
zeta symmetry; the second uses $A,C\in\mathcal O(\mathcal H)$ and
\eqref{eq:a12-ABCD}, initially on $\Re s>2$.
Separate holomorphy of both $\alpha$ and $\delta$ suffices; holomorphy of
$\delta$ alone does not.

Subtracting $\delta$ leaves the zero row and the ordinal part.
Without a new independent estimate, this returns to a classical shifted
reciprocal-zeta problem. There is no assumed global bijection between
the poles of the full $\alpha$ and the zeros of $\zeta$.

\subsubsection{8A: contours}
\label{subsubsec:a12-contours}
A contour route must establish continuation, control away from poles,
and a count excluding poles. Uniform polynomial growth on a strip
containing poles is impossible without removing their neighbourhoods
or their principal parts. Furthermore,
\begin{equation}
\label{eq:a12-argument-principle}
\frac1{2\pi i}\oint_\Gamma\frac{\alpha'}{\alpha}\,ds
=N_{\rm zeros}(\Gamma)-N_{\rm poles}(\Gamma).
\end{equation}
A zero integral is not a zero pole count. A lower bound on a separating
curve does not supply the missing interior count.
For RH it suffices to work in the open half-plane $\mathcal H$;
a strip crossing its boundary is not intrinsically necessary.

\subsubsection{8B: positivity and the correctly shifted Li map}
\label{subsubsec:a12-Li}
One first needs a completed function with proved divisor, controlled
growth, and an appropriate convergent product.
For candidate poles $\varpi=1+\rho$, the transferred expression is
\begin{equation}
\label{eq:a12-Li}
\lambda_n^\alpha=
\lim_{T\to\infty}\sum_{\substack{\varpi=1+\rho\\|\Im\varpi|\le T}}
\left[1-\left(\frac{\varpi-2}{\varpi-1}\right)^n\right],
\end{equation}
with multiplicities and the displayed summation convention.
The map is not an arbitrary $1-\varpi_0/\varpi$:
\begin{equation}
\label{eq:a12-Li-circle}
\left|\frac{\varpi-2}{\varpi-1}\right|=1
\quad\Longleftrightarrow\quad\Re\varpi=\frac32.
\end{equation}
Li's theorem gives a criterion, not automatic positivity \cite{Li1997}.
An ordinal implementation must derive the relevant form and prove every
required inequality. Generic Gram positivity is insufficient.
The sign involution supplies neither the divisor nor the completion.

\subsubsection{8C: spectral realisation}
\label{subsubsec:a12-spectral}
A spectral programme would require:
\begin{enumerate}[label=(\roman*)]
\item a specified Hilbert space and densely defined self-adjoint operator;
\item a justified determinant, regularised when necessary;
\item an exact relation between its divisor and the relevant analytic divisor;
\item an affine spectral map taking real spectrum to $\Re s=3/2$.
\end{enumerate}
Reality of a spectrum matters only after the divisor relation is proved.
An arbitrary real-spectrum matrix, or an operator defined using already
known zero ordinates, does not provide that relation independently.

\subsection{Step 9: corrected rules and excluded shortcuts}
\label{subsec:a12-rules}
\begin{enumerate}[label=\textup{R\arabic*.},leftmargin=*]
\item Do not assume an RH-equivalent estimate for $M,SS,W,d_j$ or a
corresponding abscissa while trying to prove it. An independently proved
equivalent assertion is not circular; neither is an unconditional bound.
\item Distinguish original-order convergence, absolute convergence by
pairs, and analytic continuation. Retain unfinished-pair boundary terms.
\item Use the exact integral \eqref{eq:a12-exact-integral}. The presence of
a displacement in an identity is not circular; assuming an RH-sized
bound for it would be.
\item Do not transfer poles or a functional equation from a single
summand to $\alpha$. Coefficient size alone does not force a symmetry line.
\item A power-saving prime-sum remainder requires a specified main term.
A finite inverse-logarithm expansion does not supply it. This obstructs
that approximation scheme, not every possible analytic continuation.
\item A prime-index block of length $|M(x)|$ is not automatically a
physical prime interval of length $\sqrt x$. No argument here identifies
a single missing $\epsilon$ in a prime-interval bound as the sole obstacle.
\item Fixed-lag logarithmic cancellation does not estimate the absolute
height in \eqref{eq:a12-local-error}, growing lags, or the full signed
weighted tail without further proof.
\item An identity valid for every sign word is not by itself arithmetic
specificity. Auxiliary words test a universal mechanism; they are not
asserted to be the M\"obius sequence.
\item Quadratic comparison theorems prove their stated comparisons only.
Proposition~\ref{prop:a12-order-example} refutes the claim that every
quadratic functional loses order.
\item The off-boundary assertion
\begin{equation}
\label{eq:a12-L2-correction}
\int_1^\infty M(x)^2x^{-2\sigma-1}\,dx<\infty
\quad\text{for every }\sigma>1/2
\end{equation}
is equivalent to RH without simplicity of zeros. RH implies finiteness
by choosing a smaller exponent. Conversely, Cauchy--Schwarz makes the
Mellin integral of $M$ analytic to the right of $1/2$.
The stronger weak Mertens condition
$\int_1^X M(x)^2x^{-2}\,dx=O(\log X)$ has additional simplicity and
negative-moment consequences \cite{Ng2004}.
The active route here does not assume a quadratic estimate.
\item Balanced finite words justify combinatorial identities, not
replacement of all integer cutoffs by balanced cutoffs in an infinite
analytic argument.
\end{enumerate}

\subsection{Step 10: comparison with established research programmes}
\label{subsec:a12-literature}
Each comparison separates what the cited work proves from what would
still be needed for RH. The missing implication identified here is a
logical assessment, not a claim that the programme is impossible or
that its authors' valid theorem is erroneous.

\subsubsection{10A: Mertens criteria and conditional moments}
\label{subsubsec:a12-lit-Mertens}
The Mellin identity \eqref{eq:a12-Mellin} and the Littlewood criterion
connect cancellation of $\mu$ to zero exclusion
\cite{Littlewood1912,Titchmarsh1986}.
Replacing a Mertens bound by one for $SS$ or $W$ does not independently
prove it. The additional occurrence-pair representation is the distinction
here, not a weaker convergence threshold.

Ng's distributional and moment theorems explicitly use RH and additional
negative-moment hypotheses \cite{Ng2004}. Using their conclusions without
those hypotheses to prove RH would reverse the logical direction.
Here finite pairing identities and bounded-duration holomorphy are
unconditional, but the infinite-tail estimate remains missing.

\subsubsection{10B: Nyman--Beurling and B\'aez-Duarte}
\label{subsubsec:a12-lit-NB}
Let $\{y\}=y-\lfloor y\rfloor$. In $L^2(0,\infty)$ set
\begin{equation}
\label{eq:a12-NB}
\rho_a(x)=\left\{\frac1{ax}\right\},\qquad
\mathcal B_{\mathbb N}=\operatorname{span}\{\rho_a:a\in\mathbb N\}.
\end{equation}
B\'aez-Duarte's strengthening makes RH equivalent to
$\mathbf1_{(0,1)}\in\overline{\mathcal B_{\mathbb N}}$
\cite{BaezDuarte2002}. The missing RH step is the approximation, not
nonnegativity of a squared norm.

The present route has two actual occurrence clocks and fixed prime
weights rather than freely selected approximating coefficients.
Its target is nevertheless an infinite convergence statement.
An arithmetic definition without mentioning zeta zeros is not unique
to this construction, nor does it remove that obligation.

\subsubsection{10C: Li--Weil positivity}
\label{subsubsec:a12-lit-Li}
Bombieri and Lagarias relate Li's criterion to general inequalities for
complex multisets and to the Guinand--Weil explicit formula
\cite{BombieriLagarias1999}. A positivity reformulation need not be
specific to the original zeta notation. The unsupplied step is positivity
for every required index or admissible test.

Ordinal pairing gives explicit finite kernels, but their automatic Gram
nonnegativity is not the needed arithmetic positivity. No equality between
an ordinal form and the required Li--Weil form is established here.
The shift $\rho\mapsto1+\rho$ in \eqref{eq:a12-Li} is an exact coordinate
change, not a new positivity mechanism.

\subsubsection{10D: Hilbert--P\'olya, Berry--Keating, and Connes}
\label{subsubsec:a12-lit-spectral}
Berry and Keating study an $xp$-based spectral interpretation and
explicitly identify substantial speculative gaps \cite{BerryKeating1999}.
Matching a spectral density or a semiclassical count would not identify
every eigenvalue with a zeta zero. The exact divisor relation is essential.

Connes supplies a different framework: critical zeros appear as an
absorption spectrum, and hypothetical off-critical zeros as resonances.
His work reduces RH to a trace-formula assertion \cite{Connes1998}.
The existence of that spectral interpretation does not already exclude
the resonances.

Here the starting data are discrete occurrence lists and an exact finite
summation identity, without postulating an operator. This changes the
starting object but supplies no self-adjoint realisation. Adopting a
spectral route would still require all four obligations of Step 8C.

\subsubsection{10E: de Bruijn--Newman deformation}
\label{subsubsec:a12-lit-Newman}
In the de Bruijn--Newman family all zeros are real precisely beyond a
threshold $\Lambda$. RH is equivalent to $\Lambda\le0$.
Rodgers and Tao prove $\Lambda\ge0$ \cite{RodgersTao2018}.
That proves Newman's conjecture, not the complementary inequality needed
for RH; it is not a defect in their theorem.

The present duration cutoff $L$ is not a heat-flow parameter.
Holomorphy for every fixed $L$ does not give monotonic zero motion or a
threshold theorem. Any such additional structure would need a proof.

\subsubsection{10F: local correlations and short intervals}
\label{subsubsec:a12-lit-local}
Tao's logarithmic two-point theorem and the short-interval results of
Matom\"aki and Radziwi\l\l establish genuine arithmetic cancellation in
their respective regimes \cite{Tao2016,MatomakiRadziwill2016}.
Neither states the weighted original-order Cauchy estimate
\eqref{eq:a12-target}. Equation \eqref{eq:a12-local-error} identifies
a specific unestimated height term, so describing the observable as local
does not license an exchange of limits.

The distinctive interface is an exact relation between matched-pair tails
and two-time heights. The inherited theorem controls its increment error,
not automatically its absolute height, sign, rank weight, or full tail.

\subsubsection{10G: rigorous computation and statistical models}
\label{subsubsec:a12-lit-numerics}
Platt and Trudgian rigorously verify RH through height $3\cdot10^{12}$
\cite{PlattTrudgian2021}. Their theorem certifies that finite interval; it
does not quantify over all heights.

Likewise, the present finite-word checks validate finite identities, and
the numerical controls test mechanisms over a bounded range. Neither
bounds the supremum over all $Y\ge X$ in \eqref{eq:a12-target}.
A statistical resemblance, or its absence, does not supply a
deterministic infinite-tail estimate.

\subsection{Step 11: precise distinctiveness and limits of novelty}
\label{subsec:a12-distinctiveness}
The combination studied here has six demonstrated features:
\begin{enumerate}[label=(\roman*),leftmargin=*]
\item Ordinary primes are attached to separately numbered positive and
negative occurrences of $\mu$, with an explicit zero row.
\item Position, ordinal growth, and prime displacement separate exactly:
the ordinal partial sum is $M(Q+1)/2$, and the excess is
$D_w(N_+)-D_w(N_-)$.
\item Integer-clock, squarefree-clock, inverse-pair, and active-cut
descriptions have exact transition identities.
\item Nonpolynomial duration observables distinguish paths having the
same one-time height histogram.
\item Every fixed-duration paired subseries is holomorphic for $\Re s>1$,
with an explicit cutoff-dependent majorant.
\item The original-order boundary term and the two-time height remainder
remain visible instead of being suppressed by a rearrangement.
\end{enumerate}

These features describe this construction; they do not establish priority
for each ingredient. Monotone transport, layer counting, Mellin transforms,
partial summation, convergence-abscissa formulae, and positivity criteria
have classical antecedents. The sign flip is generic, and the
bounded-duration theorem applies to much more than the M\"obius word.

The complete transform is a lossless encoding, not additional independent
arithmetic data: $\mu(n)=0$ is detected by $S(n)=1$, and otherwise $\mu(n)$
is the sign of $S(n)$. At one endpoint, $SS(n)$ depends only on $(M(n),Q(n))$.
Duration data distinguish different histories, not two copies of an
identical full input trajectory.

The defensible difference from the cited programmes is the particular
rank-prime encoding and its exact pair/cut interfaces.
No exhaustive literature theorem establishes a difference from every other
construction, and no theorem here proves the remaining RH-equivalent
estimate easier than the classical one. Its potential value would be a
new, independently controlled arithmetic mechanism at those interfaces;
that mechanism is still unresolved.

\subsection{Step 12: complete execution order and present status}
\label{subsec:a12-execution}
\begin{enumerate}[label=\textup{P\arabic*.},leftmargin=*]
\item Fix clocks, weights, and finite-cut conventions: completed.
\item Establish the initial series and $A+B+C+\delta$: completed.
\item Separate poles of $B$ from unproved claims about $\alpha$: completed.
\item Prove the exact pair integral and unit-cut reorganisation: completed.
\item Prove the duration identity and order-sensitive example: completed.
\item Determine the scope of the fixed-lag theorem: completed; an
absolute-height remainder is left unestimated.
\item Prove fixed-duration holomorphy: completed.
\item Exclude universal tail deletion by sign antisymmetry: completed;
M\"obius-specific cancellation is not excluded.
\item Establish weighted signed cancellation uniformly over durations,
including unfinished pairs: unresolved.
\item Deduce \eqref{eq:a12-target}: unresolved because P9 is missing.
\item Apply Theorem~\ref{thm:a12-transfer} only after proving its premise:
the implication is proved; the premise is not.
\item Alternatively prove \eqref{eq:a12-analytic-criterion}. Contours,
positivity, and spectral constructions are admissible only with their
additional obligations verified; none is a completed substitute.
\end{enumerate}

The next mathematical work is P9, not declaring a pole line or promoting a
finite computation into an infinite bound. This chapter is an audited
research programme with explicit intermediate results, not a proof of RH.

\subsection{Step 13: the energy of SS and its exact critical threshold}
\label{subsec:a12-SS-energy}
For real $x$, the correct step-function subtraction is
\begin{equation}
\label{eq:a12-T-real}
T(x)=SS(x)-\lfloor x\rfloor+Q(x)
=\mathcal P(N_+(x))-\mathcal P(N_-(x)),\qquad
\mathcal P(m)=\sum_{j\le m}p_j.
\end{equation}
Define the nonnegative, possibly infinite quantity
\begin{equation}
\label{eq:a12-I}
\boxed{\mathcal I(\sigma)=\int_1^\infty
\bigl(SS(x)-\lfloor x\rfloor+Q(x)\bigr)^2x^{-2\sigma-1}\,dx.}
\end{equation}
Using $x$ instead of $\lfloor x\rfloor$ changes $T$ by $-\{x\}$.
Since $\{x\}$ is bounded, this change does not affect finiteness of
\eqref{eq:a12-I} for $\sigma>0$.

\begin{proposition}[Energy comparison, including the logarithmic weight]
\label{prop:a12-energy-comparison}
For sufficiently large $x_0>1$ and any real $\sigma$,
\begin{equation}
\label{eq:a12-energy-comparison}
\mathcal I(\sigma)<\infty
\quad\Longleftrightarrow\quad
\int_{x_0}^\infty M(x)^2(\log x)^2
x^{-2(\sigma-1)-1}\,dx<\infty.
\end{equation}
More precisely, the two nonnegative integrands have asymptotic ratio
$c_0^2=9/\pi^4$, wherever $M(x)\ne0$, and both vanish when $M(x)=0$.
\end{proposition}
\begin{proof}
The exact prime block in \eqref{eq:a12-T-real} has $|M(x)|$ terms.
PNT gives $M(x)=o(x)$, $Q(x)\sim2c_0x$, and
$p_j\sim j\log j$. Every index in the block is therefore
$(1+o(1))c_0x$, uniformly over that block. Consequently
\begin{equation}
\label{eq:a12-T-asymptotic}
T(x)=c_0x\log x\,M(x)(1+o(1)).
\end{equation}
This is a two-sided comparison in absolute value for all sufficiently
large $x$, with the zero case treated exactly. Squaring and integrating
proves the assertion. The finite initial interval causes no issue.
\end{proof}

\begin{theorem}[A zero forces energy divergence even at its endpoint]
\label{thm:a12-energy-zero}
If $\zeta(\rho)=0$ for a nontrivial zero $\rho=\beta+i\gamma$, then
\begin{equation}
\label{eq:a12-energy-zero}
\boxed{\mathcal I(\sigma)=\infty\qquad(\sigma\le1+\beta).}
\end{equation}
In particular, unconditionally $\mathcal I(3/2)=\infty$.
\end{theorem}
\begin{proof}
Suppose instead that $\mathcal I(1+\beta)<\infty$. By
Proposition~\ref{prop:a12-energy-comparison},
\begin{equation}
\label{eq:a12-Jbeta}
J_\beta:=\int_{x_0}^\infty
M(x)^2(\log x)^2x^{-2\beta-1}\,dx<\infty.
\end{equation}
For $\varepsilon\ge0$, Cauchy--Schwarz gives the uniform estimate
\begin{align}
\left|\int_{x_0}^\infty
M(x)x^{-\beta-\varepsilon-i\gamma-1}\,dx\right|
&\le J_\beta^{1/2}
\left(\int_{x_0}^\infty
\frac{x^{-2\varepsilon-1}}{(\log x)^2}\,dx\right)^{1/2}
\notag\\
&\le \frac{J_\beta^{1/2}}{\sqrt{\log x_0}}.
\label{eq:a12-energy-boundary-CS}
\end{align}
The same bound with absolute values inside the first integral proves
absolute convergence at the boundary. The full integral
$F(z)=\int_1^\infty M(x)x^{-z-1}\,dx$ is holomorphic for $\Re z>\beta$
and remains bounded as $z=\rho+\varepsilon$ tends to $\rho$ from the right.
On $\Re z>1$, it equals $1/(z\zeta(z))$. The identity theorem, applied to
$z\zeta(z)F(z)=1$ away from $z=1$, propagates that identity throughout
$\Re z>\beta$. It either already excludes a zero in this open half-plane,
or contradicts boundedness when $z\to\rho$.
Thus the endpoint integral diverges.
For $\sigma\le1+\beta$, its integrand is at least that of the endpoint
because $x\ge1$. Existence of a critical-line zero gives the final claim.
\end{proof}

\begin{corollary}[RH expressed as the energy of SS]
\label{cor:a12-energy-RH}
\begin{equation}
\label{eq:a12-energy-RH}
\boxed{\mathrm{RH}\quad\Longleftrightarrow\quad
\mathcal I(\sigma)<\infty\ \hbox{for every }\sigma>3/2.}
\end{equation}
More generally,
\begin{equation}
\label{eq:a12-energy-abscissa}
\inf\{\sigma\in\mathbb R:\mathcal I(\sigma)<\infty\}=1+\Theta.
\end{equation}
\end{corollary}
\begin{proof}
Under RH, $M(x)\ll_\epsilon x^{1/2+\epsilon}$. For a given $\sigma>3/2$,
choose $\epsilon<\sigma-3/2$ and use
\eqref{eq:a12-energy-comparison}. Conversely, a zero with $\beta>1/2$
contradicts the claimed finiteness at any
$3/2<\sigma\le1+\beta$, by Theorem~\ref{thm:a12-energy-zero}.
The functional equation then gives RH. The same argument with
$M(x)\ll_\epsilon x^{\Theta+\epsilon}$ proves
\eqref{eq:a12-energy-abscissa}. If the supremum $\Theta$ is attained by a
zero, the integral diverges at $1+\Theta$; no endpoint assertion is
inferred from a supremum that is not known to be attained.
\end{proof}

This proof transfers through $M$ using an explicit two-sided comparison;
it does not assume a pole correspondence for the Mellin transform of $T$.
Indeed the exact initial-half-plane identity is
\begin{equation}
\label{eq:a12-T-Mellin}
s\int_1^\infty T(x)x^{-s-1}\,dx
=\alpha(s)-A(s)=B(s)+C(s)+\delta(s).
\end{equation}
The term $\delta$ cannot be silently included in a holomorphic remainder.
The logarithmic factor is also useful:
it makes the boundary Cauchy--Schwarz integral
$\int dx/(x\log^2x)$ finite in
\eqref{eq:a12-energy-boundary-CS}.
Off-boundary finiteness for every $\sigma>3/2$ does not impose simplicity
of zeta zeros; it is different from a bounded critical mean-square
condition.

\subsection{Step 14: explicit poles and an attempted isolation of delta's singularities}
\label{subsec:a12-pole-attempt}
This step separates the meromorphic blocks from singularity claims that
still depend on the remainder.

\subsubsection{14A: what the two explicit zeta blocks actually imply}
\label{subsubsec:a12-A-poles}
The pole statement for $B$ is Proposition~\ref{prop:a12-poles}.
For $A$, the pole at $s=1$ has residue $1-1/\zeta(2)$.
At $s_0=\rho/2$, where $\rho$ is a nontrivial zero of multiplicity $m$,
the potential pole order is
\begin{equation}
\label{eq:a12-A-order}
\max\{m-\operatorname{ord}_{s=\rho/2}\zeta(s),0\}.
\end{equation}
For a simple zero with $\zeta(\rho/2)\ne0$ its residue is
\begin{equation}
\label{eq:a12-A-residue}
\operatorname*{Res}_{s=\rho/2}A(s)
=-\frac{\zeta(\rho/2)}{2\zeta'(\rho)}.
\end{equation}
At negative odd integers $A$ has simple poles arising from $\zeta(2s)$;
at negative even integers the numerator's simple zero cancels the
denominator's simple zero. Both $s=0$ and $s=1/2$ are regular.
These are poles of the displayed block, not assertions about their
survival in $\alpha$ outside the established common domain.
In particular, claiming that all poles lie on $3/2$ would already require
specifying nontrivial poles and proving what happens to the other blocks.

\subsubsection{14B: a differentiated ordinal block}
\label{subsubsec:a12-delta-model}
The partial sum of the coefficients of $\delta$ is exactly
\begin{equation}
\label{eq:a12-W-pole}
W(x)=D_w(N_+(x))-D_w(N_-(x))
=T(x)-\frac{M(x)(Q(x)+1)}2.
\end{equation}
For $\Re s>2$, introduce
\begin{equation}
\label{eq:a12-Pdelta}
F_\zeta(z)=\frac1{z\zeta(z)},\qquad
P_\delta(s)=-c_0sF_\zeta'(s-1)=\frac{B(s)}s-B'(s).
\end{equation}
Differentiation under the absolutely convergent Mellin integral gives the
exact decomposition
\begin{align}
\delta(s)&=P_\delta(s)+E_\delta(s),\label{eq:a12-delta-split}\\
E_\delta(s)&=s\int_1^\infty
\bigl(W(x)-c_0x\log x\,M(x)\bigr)x^{-s-1}\,dx.
\label{eq:a12-Edelta}
\end{align}
The factor $B(s)/s$ has its natural removable interpretation at $s=0$.

\begin{proposition}[Explicit higher-order pole candidates]
\label{prop:a12-Pdelta}
At $s_0=1+\rho$, if $\rho$ is a zero of multiplicity $m$, the explicit
meromorphic function $P_\delta$ has a pole of order $m+1$.
For a simple zero, with
$r_\rho=c_0(1+\rho)/(\rho\zeta'(\rho))$,
\begin{equation}
\label{eq:a12-Pdelta-Laurent}
P_\delta(s)=
\frac{r_\rho}{(s-s_0)^2}
+\frac{r_\rho/s_0}{s-s_0}+O(1).
\end{equation}
\end{proposition}
\begin{proof}
Insert $B(s)=r_\rho/(s-s_0)+b_0+O(s-s_0)$ into
$B(s)/s-B'(s)$. In higher multiplicity, differentiation raises the
leading pole order by one with a nonzero leading coefficient.
\end{proof}

This locates poles of $P_\delta$, and supplies a precise candidate
principal part for $\delta$. It does not establish any pole of $\delta$.
For example, holomorphy of $E_\delta$ in a neighbourhood of $s_0$ would
make \eqref{eq:a12-Pdelta-Laurent} the principal part of $\delta$ there.
No such neighbourhood is presently proved. The comparison
$W(x)=c_0x\log x\,M(x)(1+o(1))$ alone gives no power saving in the
Mellin remainder. The integral \eqref{eq:a12-Edelta} is justified for
$\Re s>1+\Theta$ by the classical growth bound, but this does not cross
the candidate singularities.

For orientation, the first five positive critical ordinates give the
following numerical coefficients. These are evaluations of the explicit
blocks, not a computation of a continuation of $\delta$.
\begin{center}
\small
\begin{tabular}{rrr}
\toprule
$\gamma$ in $\rho=1/2+i\gamma$&
$\Re r_\rho$&$\Im r_\rho$\\
\midrule
14.1347251417&0.3751535752&$-0.0871438181$\\
21.0220396388&0.2639742227&$ 0.0461618613$\\
25.0108575801&0.2065891363&$-0.0811247572$\\
30.4248761259&0.2042801902&$ 0.1128180994$\\
32.9350615877&0.1811319824&$-0.1250906057$\\
\bottomrule
\end{tabular}
\end{center}
At each listed $s_0=3/2+i\gamma$, $r_\rho$ is the residue of $B$ and the
double-pole coefficient of $P_\delta$.
The accompanying code evaluates the Laurent limits independently.

\subsubsection{14C: why a relative asymptotic does not prove meromorphy}
\label{subsubsec:a12-branch-model}
Here is an explicit control example, unrelated to any assumed formula
for the true $W$. For $x\ge e$ take
$W_{\rm mod}(x)=x^{s_0}(\log x+\log\log x)$.
The second term is $o(x^{s_0}\log x)$ in modulus. Nevertheless, putting
$z=s-s_0$ and $\Re z>0$ gives
\begin{align}
\int_e^\infty W_{\rm mod}(x)x^{-s-1}\,dx
&=\int_1^\infty e^{-zu}(u+\log u)\,du\notag\\
&=\frac{e^{-z}(1+z)}{z^2}+\frac{E_1(z)}z,
\label{eq:a12-branch-example}
\end{align}
where the exponential integral satisfies
$E_1(z)=-\gamma_{\!E}-\log z+O(z)$ near zero in a slit domain
\cite{DLMF}. Thus a double-pole main term can coexist with a logarithmic
branch singularity in a relatively smaller remainder.
An actual meromorphic continuation for $\delta$ must exclude such
possibilities by additional information, rather than by interpreting
an asymptotic logarithmic weight as an exact pole order.

\subsection{Step 15: audit of five routes suggested by the Connes survey}
\label{subsec:a12-Connes-survey}
The source for this comparison is Connes's 2026 survey, specifically
Sections 2.1.3, 2.2.3--2.2.4, 4.1, and 6.4--6.6
\cite{Connes2026}. Its conditional programmes and proved equivalences
are not failed theorems. The following distinctions determine which
parts transfer to the present object.

\subsubsection{15A: stability of poles is local and conditional}
\label{subsubsec:a12-pole-stability}
If $F$ is meromorphic near $s_0$ and $H$ is holomorphic there, the negative
powers of the Laurent series of $F+H$ equal those of $F$.
This elementary fact proves stability of every pole and its principal
part under a holomorphic additive perturbation.
The Borchsenius--Jessen result discussed in the survey illustrates very
different behaviour for zeros after adding a nonzero constant.
It does not prove that every conceivable perturbation destroys a zero
alignment.

For $\alpha=B+(A+C+\delta)$, the stability argument is available precisely
where $A+C+\delta$ is known to be holomorphic.
It supplies no holomorphy of that perturbation and no unconditional
bijection from the zeros of $\zeta$ to the poles of $\alpha$.
The exact pole bijection currently belongs to $B$.

\subsubsection{15B: Voronin universality does not forbid a polynomial lower bound}
\label{subsubsec:a12-universality}
Voronin's theorem gives approximation to fixed nonvanishing targets on
suitable compact subsets of $1/2<\Re z<1$ by vertical translates of
$\zeta$. The approximation tolerance and the translation height are
different quantifiers. Qualitative universality does not assert an
error smaller than every inverse power of that height.

There is a concrete countercheck to the proposed obstruction. Euler
summation gives
\begin{equation}
\label{eq:a12-zeta-elementary}
\zeta(z)=\frac{z}{z-1}
-z\int_1^\infty \{x\}x^{-z-1}\,dx\qquad(\Re z>0,\ z\ne1).
\end{equation}
On $3/2+\eta\le\Re s\le2-\eta$, for $0<\eta<1/4$ and $|\Im s|\ge1$,
it implies $|\zeta(s-1)|\le C_\eta(1+|\Im s|)$. Hence at every regular
point of $B$ in this closed substrip,
\begin{equation}
\label{eq:a12-B-lower}
\boxed{|B(s)|\ge\frac{c_\eta}{1+|\Im s|}.}
\end{equation}
Thus the factor containing $1/\zeta(s-1)$ actually has a polynomial
lower bound compatible with universality. This proves nothing comparable
for the sum $\alpha$, where cancellation between summands is possible.
It does disprove the assertion that universality alone closes the entire
contour route. The other contour obligations of this chapter remain.

\subsubsection{15C: a finite prime window is not yet Weil positivity}
\label{subsubsec:a12-Weil-window}
Use the symmetric multiplicative convention of the survey's Section 4.1:
$g^*(x)=\overline{g(x^{-1})}$, Haar measure $dx/x$, and Mellin--Fourier
transform $\widehat g(t)=\int_0^\infty g(x)x^{-it}\,dx/x$.
The Weil criterion has the form
\begin{equation}
\label{eq:a12-Weil}
\sum_v W_v(g*g^*)\le0
\quad\hbox{for all admissible }g,\qquad
\widehat g(i/2)=\widehat g(-i/2)=0,
\end{equation}
with the survey's sign convention. Its nonarchimedean terms are
\begin{equation}
\label{eq:a12-Weil-primes}
W_p(f)=(\log p)\sum_{m\ge1}p^{-m/2}
\bigl(f(p^m)+f(p^{-m})\bigr).
\end{equation}
If $g$ is supported in $[\lambda^{-1},\lambda]$, then $g*g^*$ is
supported in $[\lambda^{-2},\lambda^2]$; only prime powers
$p^m\le\lambda^2$ occur. The archimedean term is still essential, and
RH requires all windows and all admissible tests.

By contrast, $E_{\rm mid}=\operatorname{sgn}(M)\sum_i \tau_i g_i$ is
linear in the consecutive prime gaps $g_i$. Its square is a rank-one
quadratic expression in that gap vector. It is not the Weil form merely
because only finitely many primes occur. Here $\tau_i$ denotes the triangular
gap weight of the residual, distinct from $w_j=p_j-j$.
Even asking whether a finite vector belongs to the domain of
$QW_\lambda$ requires an embedding into the function space.
A quadratic form is not itself a subspace to which a vector can be
orthogonal.

A well-posed finite test would specify a linear map
\begin{equation}
\label{eq:a12-Weil-embedding}
\mathcal J_\lambda:\mathbb R^k\longrightarrow
\{g\in C_c^\infty((\lambda^{-1},\lambda)):
\widehat g(\pm i/2)=0\}
\end{equation}
and calculate the pulled-back matrix
$H_{ij}=QW_\lambda(\mathcal J_\lambda e_i,\mathcal J_\lambda e_j)$.
Then $\tau^\mathsf{T}H\tau$ is defined and checkable.
No canonical embedding, equality with $E_{\rm mid}$, or implication
covering all admissible tests is established by the finite-window
observation. That is the precise missing identification.

\subsubsection{15D: an exact determinant with entries zero or signed primes}
\label{subsubsec:a12-determinant}
The proposed finite algebraic construction is possible.
First recall why the classical Redheffer determinant works. Let
$Z_{ij}=\mathbf1_{i\mid j}$ for $1\le i,j\le n$.
It is upper triangular with determinant $1$, and
$(Z^{-1})_{ij}=\mu(j/i)$ when $i\mid j$, zero otherwise.
Replacing its first column by the all-ones vector gives
$R_n=Z+(\mathbf1-e_1)e_1^\mathsf{T}$. The determinant lemma yields
\begin{equation}
\label{eq:a12-Redheffer}
\det R_n=1+e_1^\mathsf{T}Z^{-1}(\mathbf1-e_1)
=\sum_{j\le n}\mu(j)=M(n).
\end{equation}
This is the classical divisor-incidence construction recalled in
\cite{Connes2026}.

For the present transform, a different elementary encoding suffices.
Put
\begin{equation}
\label{eq:a12-prime-block-matrix}
U=\begin{pmatrix}2&3\\3&5\end{pmatrix},\qquad
U^{-1}=\begin{pmatrix}5&-3\\-3&2\end{pmatrix},\qquad \det U=1.
\end{equation}
Given column blocks $u_i,v_i\in\mathbb Z^2$, form
\begin{equation}
\label{eq:a12-bordered-matrix}
\mathcal R(u,v)=
\begin{pmatrix}
U&0&\cdots&0&u_1\\
0&U&\cdots&0&u_2\\
\vdots&\vdots&\ddots&\vdots&\vdots\\
0&0&\cdots&U&u_m\\
v_1^\mathsf{T}&v_2^\mathsf{T}&\cdots&v_m^\mathsf{T}&0
\end{pmatrix}.
\end{equation}
\begin{theorem}[Prime-entry determinant representation]
\label{thm:a12-determinants}
There are explicit integer matrices $\mathcal R_T(n)$ of size
$2Q(n)+1$ and $\mathcal R_{SS}(n)$ of size $2n+1$, all of whose entries
belong to $\{0\}\cup\{\pm p_j:j\ge1\}$, such that
\begin{equation}
\label{eq:a12-det-identities}
\boxed{\det\mathcal R_T(n)=T(n),\qquad
\det\mathcal R_{SS}(n)=SS(n).}
\end{equation}
\end{theorem}
\begin{proof}
The Schur complement of the diagonal block
$D=\operatorname{diag}(U,\ldots,U)$ gives
\begin{equation}
\label{eq:a12-Schur}
\det\mathcal R(u,v)=-\sum_{i=1}^m v_i^\mathsf{T}U^{-1}u_i.
\end{equation}
For any zero or signed prime $t_i$, choose
$u_i=(t_i,0)^\mathsf{T}$ and $v_i=(-2,-3)^\mathsf{T}$.
Then
$v_i^\mathsf{T}U^{-1}u_i=(-10+9)t_i=-t_i$.
For $\mathcal R_T(n)$ use the $Q(n)$ signed-prime summands $S(s_q)$.

For $\mathcal R_{SS}(n)$ use one block per integer $i\le n$.
If $\mu(i)\ne0$, use the preceding choice with $t_i=S(i)$.
If $\mu(i)=0$, choose instead
$u_i=(5,7)^\mathsf{T}$ and $v_i=(-2,-7)^\mathsf{T}$.
In this case $U^{-1}u_i=(4,-1)^\mathsf{T}$ and
$-v_i^\mathsf{T}U^{-1}u_i=1=S(i)$.
Every entry is zero or a signed prime, and summing the contributions
in \eqref{eq:a12-Schur} proves both identities.
\end{proof}
For example, at $n=4$ the transformed sequence is $(2,-2,-3,1)$, giving
$\det\mathcal R_T(4)=-3$ and $\det\mathcal R_{SS}(4)=-2$.

This settles the existence question, but the block construction encodes
any sum of signed primes, with a separate gadget for a unit summand.
It is not a newly proved divisor-incidence constraint.
The already established transfer gives
\begin{equation}
\label{eq:a12-det-RH}
\mathrm{RH}\quad\Longleftrightarrow\quad
|\det\mathcal R_T(n)|\ll_\epsilon n^{3/2+\epsilon}
\quad\hbox{for every }\epsilon>0.
\end{equation}
The determinant growth estimate remains unproved.
Neither the representation nor the general Hadamard determinant
inequality supplies this bound. The construction is finite and testable;
its RH-strength estimate still requires the original arithmetic cancellation.

\subsubsection{15E: what regularising the pair measure actually gives}
\label{subsubsec:a12-GW}
The Guinand--Weil formula applies to an admissible test function.
A signed atomic measure is not such a function until regularised.
The finite regularisation can be worked out exactly, with no continuation
of $\delta$ assumed.
Set $c_n=\mu(n)w_{j(n)}$ at squarefree $n$ and $c_n=0$ otherwise, and
\begin{equation}
\label{eq:a12-GW-regularise}
\nu_X=\sum_{n\le X}c_n\delta_{\log n},\qquad
h_{X,\varepsilon}=\varphi_\varepsilon*\nu_X,\qquad
f_{X,\varepsilon,s}(x)=x^{-s}h_{X,\varepsilon}(\log x),
\end{equation}
where $\varphi_\varepsilon\in C_c^\infty(\mathbb R)$.
This $\nu_X$ retains unfinished pairs; it is the logarithmic pushforward
of the original-order coefficient measure. Write
\begin{equation}
\label{eq:a12-Phi-eps}
\Phi_\varepsilon(v)=\int_{\mathbb R}\varphi_\varepsilon(u)e^{vu}\,du,
\qquad \delta_X(s)=\sum_{n\le X}c_nn^{-s}.
\end{equation}
For the unnormalised Mellin convention
$\widetilde f(z)=\int_0^\infty f(x)x^{z-1}\,dx$, direct substitution gives
\begin{equation}
\label{eq:a12-GW-Mellin}
\widetilde f_{X,\varepsilon,s}(z)
=\Phi_\varepsilon(z-s)\delta_X(s-z).
\end{equation}
Using the survey's Section 2.1.3 convention, the explicit formula becomes
\begin{align}
\sum_\rho\Phi_\varepsilon(\rho-s)\delta_X(s-\rho)
={}&\Phi_\varepsilon(-s)\delta_X(s)
+\Phi_\varepsilon(1-s)\delta_X(s-1)\notag\\
&-\sum_v W_v(f_{X,\varepsilon,s}).
\label{eq:a12-GW-finite}
\end{align}
Zeros are counted with multiplicity. At fixed $X,\varepsilon,s$ the
smooth compact support gives rapid vertical decay of the Mellin
transform, so the zero sum is legitimate. Here the unnormalised local
terms are
\begin{equation}
\label{eq:a12-GW-local}
W_p(f)=(\log p)\sum_{m\ge1}
\bigl(f(p^m)+p^{-m}f(p^{-m})\bigr);
\end{equation}
the $p^{-m/2}$ in \eqref{eq:a12-Weil-primes} belongs to the different,
symmetric normalisation.

Equation \eqref{eq:a12-GW-finite} is a proved finite regularised identity.
It is not a closed continuation formula for $\delta(s)$: it also
contains $\delta_X(s-1)$, all the shifted values $\delta_X(s-\rho)$,
and local terms depending on the regularisation and cutoff.
Passing to $X\to\infty$ and then removing the regularisation requires
uniform summability or another justified limiting procedure.
That arithmetic estimate is not supplied by admissibility of each
finite test. This route remains meaningful, with the limiting step
explicit, rather than bypassing the unresolved control of $\delta$.

\subsubsection{15F: the survey's current programme and the remaining work here}
\label{subsubsec:a12-survey-status}
The survey's Sections 6.4--6.6 distinguish several established
constructions from further obligations: a simplicity and evenness
property for the lowest eigenfunction of the truncated Weil form, and
sufficient control relating that eigenfunction to another family of
approximants. Convergence for one family cannot simply be transferred
to another because they appear related. This comparison is recorded
without pursuing an operator route in the present programme.

The same logical issue occurs here: every fixed-duration
$\delta_{\le L}$ can be holomorphic in a large half-plane while the
limit as $L\to\infty$ has singularities. The new determinant identity,
the finite Guinand--Weil identity, and the energy equivalence are
established intermediate statements. None supplies the missing
uniform arithmetic estimate.

\begin{center}
\fbox{\parbox{0.91\linewidth}{
What remains unresolved is the arithmetic estimation of the signed,
weighted two-time terms, including unfinished pairs.
The two-cut identity is exact; by itself it does not control $\delta$.
The explicit poles of $B$ and $P_\delta$ have been located.
The singularities of the full $\delta$, and their effect on $\alpha$,
remain unproved without control of $E_\delta$ or another valid
continuation mechanism.}}
\end{center}
